\secciongrande{CAPACIDADES Y LIMITACIONES DE EJECUCIÓN: ¿Por qué JavaScript en el cliente?}

\resumen{
  JavaScript se utiliza en el lado cliente para mejorar la usabilidad, permitiendo interactividad sin acceder al sistema de archivos del servidor. Es crucial asegurar la navegabilidad sin JS y procesar datos recibidos del servidor para actualizar dinámicamente la interfaz.
}

\subtituloA{Ventajas de JavaScript en el Cliente}

\lista{
  **Mejora de Usabilidad:** Proporciona navegación rápida, moderna y dinámica. |
  **Preparación de Datos:** Prepara datos para su envío al servidor sin modificar archivos en el entorno del servidor. |
  **Resiliencia:** El sitio debe ser navegable incluso si JavaScript está deshabilitado. |
  **Actualización Dinámica:** Examina y actualiza secciones de la página con datos obtenidos del lado del servidor. |
  **Funcionalidades Adicionales:** Ofrece feedback al usuario, interfaz amigable, procesamiento de formularios y modificación dinámica de estilos y contenidos.
}
